\chapter{Cơ chế Hoạt động của MongoDB Sharding}
\label{chap:sharding-mechanism}

Chương này trình bày có hệ thống nền tảng lý thuyết và cơ chế vận hành bên trong của phân mảnh ngang trong \textbf{MongoDB 7.0}. Trước hết, Mục~\ref{sec:partition-vs-sharding} làm rõ ranh giới khái niệm giữa phân vùng trên một nút trong các hệ quản trị quan hệ và phân mảnh ngang đa nút. Tiếp theo, Mục~\ref{sec:sharded-arch} giới thiệu kiến trúc tổng thể cùng các thành phần cấu thành cụm MongoDB Sharded. Cuối cùng, Mục~\ref{sec:sharding-internals} giải phẫu quy trình định tuyến cho cả \texttt{find} và \texttt{aggregate}, hai chiến lược khóa phân mảnh (\textit{Hashed} và \textit{Ranged}), cùng cơ chế chia chunk và cân bằng theo dung lượng dữ liệu được cập nhật từ MongoDB 6.0.3/7.0.

\section{Bài toán Mở rộng Dữ liệu: Phân vùng Đơn nút và Phân mảnh Đa nút}
\label{sec:partition-vs-sharding}

\subsection{Phân biệt Single-node Partitioning và Multi-node Horizontal Sharding}

Khi dung lượng dữ liệu và lưu lượng truy vấn vượt quá khả năng xử lý hiệu quả của một bảng hoặc một máy chủ đơn lẻ (như nút thắt bộ nhớ đệm đã phân tích ở Chương~\ref{chap:intro-mongodb}), các hệ quản trị cơ sở dữ liệu thường áp dụng chiến lược chia nhỏ dữ liệu. Tuy nhiên, trong thực hành thiết kế, hai kỹ thuật \textit{Partitioning} (phân vùng trên một máy chủ) và \textit{Sharding} (phân mảnh phân tán trên nhiều máy chủ) thường bị sử dụng lẫn lộn dù chúng giải quyết hai lớp bài toán kiến trúc khác nhau \cite{Kleppmann2017}. Bảng~\ref{tab:partition-vs-sharding} đối chiếu trực diện hai cơ chế trên các khía cạnh kỹ thuật cốt lõi.

\begin{table}[!htb]
\centering
\caption{Đối chiếu bản chất kỹ thuật giữa Single-node Partitioning và Multi-node Horizontal Sharding.}
\label{tab:partition-vs-sharding}
\renewcommand{\arraystretch}{1.25}
\begin{tabularx}{\linewidth}{L{3.1cm} X X}
\toprule
\textbf{Khía cạnh} & \textbf{Single-node Partitioning (RDBMS)} & \textbf{Multi-node Horizontal Sharding (MongoDB)} \\
\midrule
Mục tiêu chính & Giảm không gian quét chỉ mục/bảng trên cùng một máy chủ; hỗ trợ quản trị vòng đời dữ liệu (xóa nhanh phân vùng cũ, lưu trữ phân tầng). & Vượt trần phần cứng (RAM, CPU, Disk I/O) của một máy chủ đơn lẻ bằng cách chia tải ra nhiều nút độc lập (\textit{scale-out}). \\
Phạm vi tài nguyên & Toàn bộ phân vùng dùng chung CPU, RAM, bộ nhớ đệm và kênh I/O của một nút vật lý. & Mỗi mảnh (\textit{shard}) sở hữu tiến trình \texttt{mongod}, bộ nhớ đệm WiredTiger, CPU và ổ đĩa riêng biệt. \\
Cơ chế định tuyến & Bộ tối ưu hóa truy vấn cục bộ (\textit{Query Planner}) loại bỏ phân vùng không khớp ngay trong cùng tiến trình DB. & Tiến trình định tuyến trung gian (\texttt{mongos}) đọc siêu dữ liệu cụm để chuyển tiếp yêu cầu qua mạng tới đúng shard đích. \\
Chi phí phát sinh & Chi phí lập kế hoạch thấp; không phát sinh độ trễ mạng giữa các phân vùng. & Phát sinh chi phí mạng qua tầng định tuyến, đồng bộ siêu dữ liệu và gom kết quả xuyên nút khi truy vấn không khớp khóa phân mảnh. \\
\bottomrule
\end{tabularx}
\end{table}

Như vậy, trong khi phân vùng đơn nút tối ưu hóa việc tổ chức vật lý trên cùng một máy chủ, phân mảnh ngang đa nút mở rộng đồng thời tổng dung lượng bộ nhớ đệm, năng lực tính toán CPU và băng thông I/O đĩa của hệ thống, đổi lại sự phức tạp trong định tuyến phân tán và cân bằng tải.

\subsection{Tối ưu phạm vi truy vấn phân tán: Cơ chế Shard Pruning}

Một lợi ích quan trọng khi chia nhỏ dữ liệu là tránh phải quét những phần dữ liệu không liên quan đến điều kiện truy vấn. Trong các hệ quản trị quan hệ có phân vùng bảng, cơ chế này được gọi là tỉa phân vùng (\textit{Partition Pruning}). Khi chuyển sang kiến trúc phân mảnh ngang đa nút của MongoDB, nguyên lý tối ưu tương ứng là \textbf{Shard Pruning} (định tuyến mục tiêu \textemdash  \textit{Targeted Routing}) \cite{MongoDBExplain}, phân chia truy vấn đọc thành hai chế độ thực thi cơ bản:

\begin{itemize}
  \item \textbf{Truy vấn định tuyến mục tiêu (Targeted / 1-shard):} Khi điều kiện lọc chứa đầy đủ khóa phân mảnh (hoặc tiền tố phù hợp của khóa phân mảnh theo khoảng), \texttt{mongos} chỉ gửi truy vấn tới đúng các shard chứa dải khóa khớp. Lưu ý quan trọng rằng việc tỉa mảnh hoạt động ở \textbf{cấp độ shard chứ không phải cấp độ chunk}: một truy vấn khoảng có thể đi qua nhiều chunk liên tiếp nhưng nếu toàn bộ các chunk đó cùng cư trú trên một shard duy nhất, truy vấn vẫn chỉ nhắm vào đúng $1$ shard. Trong kết quả \texttt{explain("executionStats")}, truy vấn \texttt{find} hiển thị giai đoạn gốc \texttt{SINGLE\_SHARD}, còn truy vấn \texttt{aggregate} hiển thị đúng $1$ phần tử trong bản đồ \texttt{shards} và không cần tách pipeline trộn.
  \item \textbf{Truy vấn phát tán\textendash thu gom (Scatter\textendash Gather / Multi-shard):} Khi truy vấn không chứa khóa phân mảnh (hoặc dải giá trị vắt ngang biên phân bổ giữa hai shard), \texttt{mongos} phát tán truy vấn đến nhiều (hoặc tất cả) các shard rồi gom/trộn kết quả. Trong \texttt{explain("executionStats")}, truy vấn \texttt{find} hiển thị \texttt{SHARD\_MERGE} (hoặc \texttt{SHARD\_MERGE\_SORT} nếu có sắp xếp), trong khi truy vấn \texttt{aggregate} hiển thị cấu trúc \textbf{\texttt{splitPipeline}} (gồm phần chạy song song trên các shard \texttt{shardsPart} và phần trộn \texttt{mergerPart}) đi kèm trường \textbf{\texttt{mergeType}} chỉ rõ bước trộn cuối cùng thực thi tại \texttt{mongos} hay được ủy quyền về một shard cụ thể (\texttt{primaryShard} hoặc \texttt{anyShard}).
\end{itemize}

Tỷ lệ \textit{Shard Pruning} đo trực tiếp trên toàn bộ bộ cửa sổ thời gian cố định, sự khác biệt giữa \texttt{find} và \texttt{aggregate} (\texttt{mergeType}), cùng chênh lệch độ trễ và lưu lượng mạng giữa hai chế độ trên là trọng tâm đánh giá của \textbf{RQ3}.

\section{Kiến trúc cụm Sharded của MongoDB}
\label{sec:sharded-arch}

Khi tổng kích thước \textit{working set} và lưu lượng ghi vượt quá khả năng chịu tải của một Replica Set đơn lẻ, MongoDB tổ chức nhiều máy chủ thành một \textbf{cụm phân mảnh} (\textit{Sharded Cluster}) gồm ba thành phần phối hợp chặt chẽ \cite{MongoDBSharding}:

\begin{itemize}
  \item \textbf{Shard (\texttt{mongod --shardsvr}):} Mỗi shard là một Replica Set độc lập lưu trữ một tập con dữ liệu của toàn cụm. Toàn bộ thao tác đọc ghi sau khi được định tuyến đều thực thi trực tiếp trên bộ máy lưu trữ WiredTiger và ổ đĩa riêng của các shard liên quan.
  \item \textbf{Config Server (\texttt{mongod --configsvr}):} Một Replica Set chuyên biệt (\texttt{cfgRS}) lưu giữ toàn bộ siêu dữ liệu (\textit{metadata}) của cụm trong cơ sở dữ liệu hệ thống \texttt{config}, bao gồm bản đồ ánh xạ dải khóa sang từng shard (\texttt{config.chunks}), danh sách collection đã bật sharding kèm mã \texttt{uuid} (\texttt{config.collections}), tham số cụm (\texttt{config.settings}), các phiên di chuyển đang chạy (\texttt{config.migrations}) và nhật ký lịch sử (\texttt{config.changelog}).
  \item \textbf{Query Router (\texttt{mongos}):} Tiến trình trung gian không lưu trữ dữ liệu người dùng (\textit{stateless router}) đóng vai trò cổng tiếp nhận mọi kết nối từ ứng dụng khách. \texttt{mongos} đọc siêu dữ liệu từ Config Server vào bộ nhớ đệm cục bộ, dựa vào điều kiện truy vấn để định tuyến tới đúng shard mục tiêu và điều phối luồng kết quả trả về cho máy khách.
\end{itemize}

Để phân phối dữ liệu của một collection ra nhiều shard, người thiết kế chỉ định một \textbf{Khóa phân mảnh} (\textit{Shard Key}) \textemdash  gồm một hoặc nhiều trường tồn tại trong mọi tài liệu của collection đó. Dựa trên giá trị của Shard Key, MongoDB quản lý không gian khóa theo các dải liên tục gọi là \textbf{Chunk} (hay \textit{Range}) với dung lượng mục tiêu cấu hình được (trong thực nghiệm này cố định ở mức $32\text{ MB}$). Tiến trình nền \textbf{Balancer} (chạy trên nút Primary của Config Server) theo dõi độ chênh lệch dung lượng giữa các shard và điều phối di chuyển dải dữ liệu (\textit{range/chunk migration}) khi cần thiết.

\section{Cơ chế định tuyến, Chiến lược Shard Key và Cân bằng tải (MongoDB 7.0)}
\label{sec:sharding-internals}

\subsection{Quy trình định tuyến và Phân rã độ trễ truy vấn}

Khi ứng dụng gửi một truy vấn tới \texttt{mongos}, quá trình xử lý diễn ra qua ba bước: (1) tra cứu bảng định tuyến dải khóa trong bộ nhớ đệm cục bộ của \texttt{mongos} (nếu phát hiện phiên bản siêu dữ liệu lỗi thời \textemdash  \textit{stale config}, \texttt{mongos} tự động làm mới từ Config Server); (2) chuyển tiếp truy vấn tới một hoặc nhiều shard mục tiêu qua mạng nội bộ; và (3) thu gom, trộn hoặc sắp xếp luồng kết quả (tại \texttt{mongos} đối với \texttt{find}, hoặc tại điểm trộn quy định bởi \texttt{mergeType} đối với \texttt{aggregate}) rồi trả về ứng dụng khách.

Với truy vấn đơn shard, độ trễ đo phía client xấp xỉ bằng thời gian thực thi tại shard đích cộng thêm một chặng mạng nội bộ và chi phí xử lý định tuyến của \texttt{mongos} (biểu hiện rõ nhất ở mốc dữ liệu nhỏ khi cả hai cấu hình đều nằm gọn trong RAM). Ngược lại, với truy vấn xuyên $2$ shard, tổng độ trễ bị chi phối bởi nút phản hồi chậm nhất (\textit{straggler effect} \textemdash  hiệu ứng nút chậm nhất) cộng thêm chi phí truyền tải và trộn kết quả tại nút trộn (\texttt{mongos} hoặc shard hợp nhất theo \texttt{mergeType}):
\begin{equation}
  T_{\text{total}} \approx \max\big(T_{\text{shard1}},\, T_{\text{shard2}}\big) + T_{\text{network + merge}(\texttt{mergeType})}.
  \label{eq:cross-shard-decomp}
\end{equation}
Tuy nhiên, với các truy vấn phân tích cần quét khối lượng lớn tài liệu nhưng chỉ trả về rất ít bản ghi kết quả (ví dụ thống kê Top-10 mã bệnh phổ biến nhất \texttt{Q3}), việc chia đôi không gian quét và thực thi song song trên hai shard có thể làm giảm đáng kể hạng tử $\max(T_{\text{shard1}}, T_{\text{shard2}})$ so với việc một nút đơn lẻ phải tự quét toàn bộ dữ liệu \textemdash  tạo ra sự đánh đổi trực tiếp giữa \textit{Shard Pruning} (giảm số shard tham gia) và \textit{Parallel Scan} (quét song song).

\subsection{So sánh chiến lược Hashed Sharding và Ranged Sharding}

Việc lựa chọn chiến lược \textit{Shard Key} quyết định trực tiếp đến độ đồng đều phân bố dữ liệu, khả năng mở rộng tải ghi và hiệu năng truy vấn đọc của toàn cụm. Đề tài khảo sát hai chiến lược đại diện tiêu biểu nhất trong MongoDB \cite{MongoDBSharding}:

\subsubsection{Hashed Sharding (\texttt{\{PATIENT: "hashed"\}})}
MongoDB áp dụng hàm băm 64-bit lên giá trị trường khóa phân mảnh trước khi ánh xạ vào các dải chunk. Các giá trị gần nhau về mặt logic hoặc tăng liên tục theo thời gian được phân bổ giả ngẫu nhiên đều khắp không gian băm.
\begin{itemize}
  \item \textbf{Ưu điểm:} Phân bố đồng đều cả dung lượng lưu trữ lẫn lưu lượng ghi mới (\texttt{insert}) trên tất cả các shard ngay từ thời điểm khởi tạo trên collection rỗng (đặc biệt khi dùng tham số \texttt{numInitialChunks: 16} để chia sẵn dải băm đều cho các shard trước khi nạp). Mọi truy vấn theo bệnh nhân chứa điều kiện lọc bằng trên \texttt{PATIENT} đều định tuyến chính xác vào $1$ shard.
  \item \textbf{Nhược điểm:} Phá vỡ tính cục bộ theo khoảng giá trị (\textit{range locality}). Các truy vấn lọc theo khoảng thời gian (\texttt{START}) hoặc theo mã chẩn đoán buộc phải phát tán sang toàn bộ các shard trong cụm.
\end{itemize}

\subsubsection{Ranged Sharding (\texttt{\{START: 1, PATIENT: 1\}})}
Dữ liệu được phân chia vào các dải khóa theo đúng thứ tự tăng dần của thời điểm khám \texttt{START} (kèm \texttt{PATIENT} tăng tính duy nhất). Các tài liệu có \texttt{START} gần nhau sẽ nằm chung trong cùng một dải và thường cư trú trên cùng một shard.
\begin{itemize}
  \item \textbf{Ưu điểm:} Bảo toàn tính cục bộ theo khoảng thời gian, thuận lợi cho các truy vấn lọc theo khung ngày/tháng/quý. Nếu toàn bộ khoảng thời gian truy vấn nằm trong phạm vi các dải khóa thuộc cùng một shard, \texttt{mongos} thực hiện \textit{Shard Pruning} triệt để ($1/2$ shard).
  \item \textbf{Nhược điểm:} Nếu không chia sẵn dải khóa trước khi nạp lô (\textit{pre-splitting}), toàn bộ dữ liệu ban đầu đi vào Primary Shard trước khi Balancer kịp di chuyển. Đặc biệt nghiêm trọng hơn, trong tải ghi giao dịch trực tuyến thực tế với khóa tăng đơn điệu ($\texttt{START} = \texttt{now()}$), $100\%$ lưu lượng ghi mới tất yếu dồn vào dải chứa biên trên vô cực (\texttt{MaxKey}) trên \textbf{một shard duy nhất} (\textit{Write Hotspot}), làm triệt tiêu lợi ích mở rộng ghi của cụm nhiều shard.
\end{itemize}

\subsection{Cơ chế chia Chunk, Jumbo Chunk và Tiến trình Balancer trong MongoDB 7.0}

Kiến trúc quản lý chunk và cân bằng tải của MongoDB đã có thay đổi mang tính bước ngoặt kể từ \textbf{MongoDB 6.0.3} và được hoàn thiện trong \textbf{MongoDB 7.0} \cite{MongoDBSharding}, khác biệt cơ bản so với cơ chế truyền thống trước đây:

\begin{enumerate}
  \item \textbf{Loại bỏ chia chunk tự động khi ghi (\textit{Auto-splitting removed}):} Khác với các phiên bản cũ ($\le 6.0.2$) vốn tự động tách chunk mỗi khi dung lượng vượt ngưỡng cấu hình ngay lúc chèn dữ liệu, MongoDB 7.0 \textbf{không còn tự động chia chunk trong quá trình ghi}. Một dải khóa trên một shard có thể lớn hơn nhiều lần \texttt{chunkSize} mà vẫn được giữ nguyên làm một đơn vị siêu dữ liệu gọn nhẹ. Chunk chỉ được tách động ngay tại thời điểm Balancer quyết định di chuyển dữ liệu sang shard khác (\textit{split on migration}).
  \item \textbf{Cân bằng theo dung lượng dữ liệu thực tế với ngưỡng $3 \times \text{chunkSize}$:} Thay vì cân bằng số lượng chunk, Balancer của MongoDB 7.0 so sánh \textbf{tổng dung lượng dữ liệu} của collection giữa các shard và chỉ kích hoạt di chuyển khi chênh lệch giữa shard nhiều dữ liệu nhất và shard ít dữ liệu nhất vượt quá \textbf{ba lần kích thước chunk cấu hình} (tức $3 \times 32\text{ MB} = \mathbf{96\text{ MB}}$ trong thực nghiệm này). Hệ quả thực nghiệm quan trọng là sau khi Balancer dừng, một mức chênh lệch dung lượng đến $\sim 96\text{ MB}$ (khoảng $7.3\%$ trên tập \texttt{D1M} $\approx 1{,}310\text{ MB}$) là \textit{hành vi hợp lệ của ngưỡng dừng Balancer}, không đồng nghĩa với việc chiến lược khóa phân mảnh bị lỗi.
  \item \textbf{Hợp nhất chunk tự động (\textit{AutoMerger}):} MongoDB 7.0 bật mặc định tiến trình AutoMerger nhằm gộp các chunk liền kề trên cùng một shard để giảm kích thước bảng siêu dữ liệu. Do đó, việc nhìn vào số lượng chunk trên \texttt{sh.status()} không còn phản ánh đúng độ cân bằng; thước đo bắt buộc phải là dung lượng MB và số bản ghi thực tế từ \texttt{getShardDistribution()} / \texttt{collStats}.
  \item \textbf{Điều kiện phát sinh Jumbo / Indivisible Chunk:} Một dải khóa chỉ trở thành bất khả phân chia (\textit{jumbo chunk}, gắn cờ \texttt{jumbo: true} trong \texttt{config.chunks} và bị Balancer bỏ qua không di chuyển) khi tồn tại \textbf{một giá trị Shard Key đơn lẻ} mang tổng dung lượng tài liệu vượt quá ngưỡng \texttt{chunkSize} ($32\text{ MB}$). Nếu mọi giá trị khóa riêng biệt đều nhỏ hơn $32\text{ MB}$, ngay cả khi dữ liệu bị lệch mạnh giữa các nhóm khóa, cụm vẫn có đủ điểm chia hợp lệ và không phát sinh jumbo chunk.
\end{enumerate}

Do hoạt động di chuyển dữ liệu của Balancer tiêu tốn đáng kể CPU, mạng và I/O đĩa, trước mọi bài đo độ trễ đọc và ghi OLTP ở Chương~\ref{chap:methodology}, hệ thống bắt buộc tắt Balancer (\texttt{sh.stopBalancer()}) và xác minh đúng chuẩn MongoDB 7.0 qua lệnh \texttt{admin.command("balancerStatus")} rằng \texttt{mode == "off"} (tương đương \texttt{sh.getBalancerState() == false}), \texttt{inBalancerRound == false} và tập \texttt{config.migrations} hoàn toàn rỗng (khắc phục thiếu sót nếu chỉ kiểm tra đơn lẻ hàm \texttt{sh.isBalancerRunning()}, vốn chỉ phản ánh có đang diễn ra một vòng cân bằng tại đúng thời điểm gọi hay không).
